\documentclass[11pt]{article}
\usepackage{amsmath,amssymb,amsthm}
\newtheorem{proposition}{Proposition}
\title{A conservative growing-block range for the actual transverse Wilson flow}\author{}\date{}
\begin{document}\maketitle
\paragraph{Scope and fixed-block input.}
The actual transverse fixed block flow calculation is a fixed-block statement: its relative-tree Haar map is real,
its output is the literal fine observable composed with that map, and its
corrected-score and raw-normalizer estimates follow from compact Laplace
localization. Here constants are tracked for a PARTICULAR increasing family.
Fix $N\ge2$. Take a complete four-dimensional cubical box of side at most $R$
in each direction and a full exterior collar, all boundary links flat except
the actual path $\eta(t)=\exp(t\alpha v)$, $\alpha=\sqrt{N/B}$,
$0\le t\le1$. Include every incident plaquette; fixed-only collar terms may
be included as their EXACT boundary-dependent scalar. Let $\|v\|_2\le1$ in
angular Lie-coordinate link units before the factor $\alpha$. No random
boundary typicality, torus holonomy replacement or physical-size claim is made.

Use the relative temporal axial forest: each interior vertex is connected by
successive direction-zero edges to the lower boundary face. Gauges at every
boundary vertex remain identity. Boundary vertices on other faces are not
included as free tree vertices. Every transformed tree edge is identity;
the last link to the upper face and other non-tree links remain independent
chords. The exact compact change has Haar Jacobian one, and the independent
chords have product normalized Haar. Denote their canonical coordinates by
$x$, the full quadratic incidence by $K$, and
\[
 w=-K_{JJ}^{-1}K_{JR}v,\quad E_v=v^*K_{\rm eff}v,\qquad
 \Phi_t(U)_b=e^{t\alpha w_b}U_b.                                       \tag{1}
\]
This is the original common realization of the actual transverse fixed block flow calculation, now with the specified
forest and increasing dimension. Its Haar divergence is exactly zero.

\paragraph{1. Literal polynomial geometry bounds.}
There are at most $C_NR^4$ independent real coordinates, and the incidence
operator has norm at most $C_N$, since each original plaquette has four links
and each link has bounded plaquette incidence. Each chord, completed by its
forest paths and a fixed boundary return path, bounds an explicit coordinate
strip/rectangle of at most $C R^2$ original plaquettes in the complete box
and collar. Spatial chords use a time--space strip; last temporal chords use
a return path along a lateral boundary face. Path lengths are at most $CR$.

The exact nonabelian plaquette product for such a filled loop uses conjugated
plaquette holonomies. Unitary triangle inequality is unaffected by these
conjugations. Squaring the sum, then summing over at most $CR^4$ chords,
gives the deliberately crude but uniform bound
\[
 \sum_{b\in J}d(U_b,I)^2
 \le C_N\left[R^6\sum_p s(U_p)
             +R^5\sum_{b\in R}d(\eta_b,I)^2\right].                    \tag{2}
\]
Here $d$ is the fixed group Riemannian distance and chord distance is uniformly
comparable to it on the compact group. The second term records the actual
fixed boundary return paths; it is not set to zero after moving the boundary.
The plaquette sum includes the full prescribed collar. At flat boundary,
linearizing(2) gives
\[
 K_{JJ}\ge c_NR^{-6}I,\qquad \|K\|\le C_N.                             \tag{3}
\]
The Schur variational identity then gives
$w^*K_{JJ}w\le v^*K_{RR}v\le C_N$, hence
\[
 \|w\|\le C_NR^3,\qquad \|(v,w)\|\le C_NR^3.                           \tag{4}
\]
These are coordinate-dependent upper bounds for this actual forest, not a
claim that the physical lowest mode really scales as $R^{-6}$.

\paragraph{2. A normalized compact tail bound with dimension retained.}
Take a joint angular chart of radius $\eta_R=c_NR^{-6}$ in the Euclidean sum
norm of all chord logs, with $c_N$ small enough. Equation(2), with
$\sum d(\eta_b,I)^2\le C_N/B$, bounds the original action outside this chart
below by $c_NBR^{-18}-C_N$ (reducing constants when necessary).
A ball of radius $B^{-1/2}$ in all independent angular coordinates has action
bounded by $C_N$ at these boundary data. Its actual Haar volume yields
\[
 -\log Z_B(\eta(t))\le C_NR^4\log(2+BR),
\]
up to any fixed-only scalar already retained in both sides of the normalized
comparison. Therefore the ORIGINAL normalized tail probability satisfies
\[
 q_{B,R}\le \exp[-c_NBR^{-18}+C_NR^4\log(2+BR)].                         \tag{5}
\]
All estimates are uniform in $t$. The displacement of the harmonic center is
$\alpha\|tw\|\le C_NR^3/\sqrt B$; it lies well inside the chart when
$B\ge C_NR^{18}$. No fixed-dimensional compactness gap has been hidden here.

\paragraph{3. Canonical moments without exponential dimension loss.}
In the chart write the ACTUAL log density as
$V_t(x)=S_B(x;\eta(t))-\log J_{\rm Haar}(x)$, omitting only its normalization.
Finite plaquette Taylor bounds and bounded incidence give Hessian perturbation
$C_N\eta_R$ from $K_{JJ}$. The canonical Haar Hessian contributes at most
$C_N/B$ in operator norm. With the chosen chart and large $B$, therefore
$\nabla^2V_t\ge c_NR^{-6}I=:\kappa_R I$ on its convex coordinate ball.
At its quadratic center $m=tw$,
\[
 |\nabla V_t(m)|\le C_NR^6/\sqrt B.                                    \tag{6}
\]
The local Haar correction is smaller. This bound follows from the quadratic
stationary equation and the third derivative of the actual local action,
not from replacing the local law by a Gaussian.

For completeness, radial integration by parts on this ball with vector
$(x-m)|x-m|^{2p-2}$ has nonnegative outward boundary flux, since $m$ is inside
it. Strong convexity and(6) then give, for each fixed integer $p\ge1$,
\[
 E[|x-m|^{2p}\mid\text{chart}]
 \le C_{N,p}\left[\frac{n+2p}{\kappa_R}
                  +\frac{|\nabla V_t(m)|^2}{\kappa_R^2}\right]^p
 \le C_{N,p}R^{10p},                                                   \tag{7}
\]
where $n\le C_NR^4$, provided $B\ge C_{N,p}R^{14}$.
This avoids comparing two Gaussian determinants by a factor exponential in
$n$. The compact complement contributes at most
$C_{N,p}B^pR^{4p}q_{B,R}$ to the moment of canonical distance. In the range
used below that term is smaller than(7).

\paragraph{4. Corrected-score and Hessian estimates with explicit powers.}
Let $D=\partial_t+W$, with $W$ the left-translation field in(1), and
$G_t=DS_B$. In canonical coordinates left multiplication has velocity
$w+O_N(\alpha|w||x|)$; its second coordinate derivative is
$O_N(\alpha^2|w|^2|x|)$ on this chart. These follow by differentiating the
actual local exponential product; the first commutator of $w$ with itself is
zero. All block derivatives are sums of local plaquette derivatives with
bounded incidence. The quadratic joint derivative is exactly $tE_v$,
independent of $x$. Consequently
\[
 |G_t-tE_v|\le C_N B^{-1/2}R^3(1+\|x\|^2),                             \tag{8}
\]
\[
 |DG_t-E_v|\le C_N B^{-1/2}R^6(1+\|x\|).                              \tag{9}
\]
The linear rather than quadratic growth in(9) is useful: the third derivative
of the action is $O_N(B^{-1/2})$ as a trilinear form, and the coordinate
acceleration term is bounded by
$C_N\alpha^2\|(v,w)\|^2\|x\|^2\le C_N\alpha\|(v,w)\|^2\|x\|$ inside
the chosen angular chart. The boundary log path is exactly $tv$.

Globally the original group derivatives obey
$|G_t|\le C_N\sqrt B R^5$ and $|DG_t|\le C_NR^6$ by the literal plaquette
count and(4). Combining(5),(7)--(9) with the EXACT identity
$\Gamma_B''=E DG_t-\operatorname{Var}G_t$ gives
\[
 \operatorname{Var}_tG_t\le C_N R^{26}/B+\mathcal E_{B,R},                \tag{10}
\]
\[
 |\Gamma_B''(t)-E_v|
 \le C_N[R^{11}/\sqrt B+R^{26}/B]+\mathcal E_{B,R},\quad
 \mathcal E_{B,R}=C_NBR^{10}
       e^{-c_NBR^{-18}+C_NR^4\log(2+BR)}.                               \tag{11}
\]
The exact normalized pullback affinity is at least
$1-\frac18[C_NR^{26}/B+\mathcal E_{B,R}]$ by the corrected-score path bound.
At flat boundary $\Gamma_B'(0)=0$: global conjugation makes its expected
boundary gradient an invariant vector in $\mathfrak{su}(N)$, hence zero.
Thus integrating(11) also controls the ACTUAL scalar
$\Gamma_B(1)-\frac12E_v$ by its right side. No approximate vacuum factor is
multiplied over the growing number of links.

\paragraph{5. One explicit mesoscopic range and its limits.}
For
\[
 2\le R\le B^{1/30},\qquad B\to\infty,                                \tag{12}
\]
all conditions above hold eventually, uniformly over the stated normalized
boundary directions. The tail in(11) is exponentially small, while both
polynomial errors are at most $C_N B^{-2/15}$. This proves a growing-block
version of the actual nonabelian transverse response and affinity estimate.
For a specified physical transverse temporal-face mode with
$E_v\ge c/R$, (11) even has relative error bounded by
$C[R^{12}/\sqrt B+R^{27}/B]+R\mathcal E_{B,R}$, tending to zero in(12).
This last qualification is about that actual mode, not every possible
positive eigenvalue in arbitrary boundary/tree coordinates.

On $B\asymp\log(1/a)$, the physical radius of(12) is only
$a[\log(1/a)]^{1/30}\to0$. The box has a flat exterior with the explicitly
small transverse boundary path; a random original exterior has NOT been
shown to belong to this class. The estimates cannot be tiled over a physical
volume without controlling the number of blocks, their actual normalizers,
induced couplings and boundary records. They do not cover nontrivial global
holonomy sectors. In particular the polynomial error per block is not a
summable full-lattice error. This is a genuine conservative growing local
range, not a physical-block or fullS94 closure.
\end{document}
